% !TeX root = ../main.tex
\section{Related Work}
\label{sec:related_work}

\subsection{Camera-Controllable Visual World Models}
\label{sec:rw_camera_models}

Camera-controllable visual world models generate future observations in response to viewpoint changes, enabling sequential exploration of a visual environment.
Earlier camera-conditioned generators, including MotionCtrl~\cite{wang2024motionctrlunifiedflexiblemotion}, CameraCtrl~\cite{he2025cameractrlenablingcameracontrol}, and CamCo~\cite{xu2024camcocameracontrollable3dconsistentimagetovideo}, establish pose-based control over video synthesis.
Interactive systems extend this capability to ongoing user control: Matrix-Game~2.0~\cite{he2026matrixgame20opensourcerealtime} injects keyboard and mouse actions into streaming generation, while Yume~\cite{mao2025yumeinteractiveworldgeneration} maps keyboard inputs to quantized camera motions for world exploration.
WorldPlay~\cite{sun2026worldplaylongtermgeometricconsistency} and LingBot-World~\cite{robbyantteam2026advancingopensourceworldmodels} further address scene persistence over longer interactive rollouts through mechanisms that retain historical context.
Aerial applications introduce flight-specific observations and controls, as explored by AirScape~\cite{zhao2025airscapeaerialgenerativeworld} and Aerial World Model~\cite{zhang2026aerialworldmodellonghorizon}.

Recent open systems address the inference cost and implementation complexity of long-horizon world modeling.
Efficiency-oriented models, including SANA-WM~\cite{zhu2026sanawmefficientminutescaleworld}, Wonder~\cite{xu2026wondervideoworldmodel}, and ABot-World-0~\cite{jiang2026abotworld0infiniteinteractiveworld}, reduce the computational cost of interactive generation through efficient attention, selective history access, or optimized streaming execution.
Complementing these model-level advances, minWM~\cite{zhao2026minwmfullstackopensourceframework}, BiWM~\cite{rui2026biwmadvancingopensourceinteractive}, and SolarWM~\cite{huang2026solarwmopendatascalable} provide reusable frameworks that organize camera-control adaptation, training, and inference into pipelines supporting multiple video backbones.

Alongside these advances in efficiency and system design, joint appearance and depth modeling has been explored for camera-conditioned generation.
Aether~\cite{zhu2025aethergeometricawareunifiedworld} jointly learns geometric reconstruction, video prediction, and visual planning.
DualCamCtrl~\cite{zhang2025dualcamctrldualbranchdiffusionmodel} couples parallel RGB and depth diffusion branches through semantic-guided mutual alignment, using an estimated initial depth map in its image-conditioned formulation.
IDCNet~\cite{liu2025idcnetguidedvideodiffusion} combines metric-aligned training data with joint RGB-depth generation conditioned on initial RGB-depth observations and camera trajectories.
Our focus is few-step causal RGB-depth generation for low-altitude flight from an initial RGB frame, text, and a camera trajectory.
We exchange appearance and geometry information through shared early transformer layers before modality-specific prediction, and supervise decoded depth with pixelwise log-depth and temporal-block scale losses.
This metric supervision is retained during causal adaptation, connecting joint geometry learning to sequential RGB-depth generation.

\subsection{Camera-Annotated Video Data}
\label{sec:rw_camera_data}

Camera-controlled generation depends on video data that associates observations with camera motion.
DL3DV-10K~\cite{ling2024dl3dv10klargescalescenedataset} provides diverse scene captures with reconstructed camera parameters, while Sekai~\cite{li2025sekaivideodatasetworld} combines real walking and drone videos with estimated trajectories and game videos with engine-recorded camera poses.
RealCam-Vid~\cite{zheng2025realcamvidhighresolutionvideodataset} reconstructs camera motion in dynamic scenes and aligns its scale using metric-depth predictions, while SpatialVID~\cite{wang2025spatialvidlargescalevideodataset} provides estimated poses and depth together with derived motion instructions.
Combining such sources requires reconciling their camera conventions, translation scales, and annotation procedures.

SolarWM~\cite{huang2026solarwmopendatascalable} standardizes heterogeneous sources into frame-aligned video, camera, caption, and quality records.
SANA-WM~\cite{zhu2026sanawmefficientminutescaleworld} builds metric camera annotations using temporally consistent geometry estimation and monocular metric-depth predictions.
Both pipelines accommodate existing camera or depth annotations, while their video-only branches recover metric geometry through estimation.
Our data construction combines simulation, public videos, and real UAV footage, using exact metric trajectories from simulation to calibrate scale recovery and compare pose estimators against ground truth.
The selected annotation procedure is then applied to recover camera trajectories from unannotated videos in a common metric convention.

\subsection{Autoregressive Video Generation}
\label{sec:rw_autoregressive}

Autoregressive video generation produces successive frames or temporal blocks conditioned on prior observations.
Diffusion Forcing~\cite{chen2024diffusionforcingnexttokenprediction} trains causal models with independently sampled noise levels across tokens, connecting next-token prediction with sequence diffusion.
CausVid~\cite{yin2025slowbidirectionalfastautoregressive} distills a bidirectional diffusion teacher into a few-step causal generator.
Self Forcing~\cite{huang2025selfforcingbridgingtraintest} trains on self-generated histories and applies distribution supervision, addressing the mismatch between reference histories during training and generated histories at inference.

Interactive world models adopt different training strategies to support camera-controlled autoregressive generation.
SolarWM~\cite{huang2026solarwmopendatascalable} uses bidirectional adaptation, teacher-forced autoregressive initialization, and distribution matching distillation as three successive training stages.
BiWM~\cite{rui2026biwmadvancingopensourceinteractive} retains bidirectional attention over the current chunk and visible history while generating chunks sequentially, using control fine-tuning followed by few-step distillation.
Direct Forcing~\cite{xu2026directforcingsinglestagetraining} instead combines future-frame feature alignment under reference histories with feature-space distribution matching under self-generated histories in one causal training stage.
Following this training formulation, we extend continuation alignment and rollout distribution matching to both RGB and depth while retaining decoded metric-depth supervision, adapting the joint model directly from bidirectional learning to few-step causal generation.
